\section{Summary}
This work reproduced AgentSemSteCom, an agentic, generative semantic steganographic
communication scheme, end to end on a single 6\,GB laptop GPU, using the components named in the
paper and documenting every unavoidable difference. The reproduction is in the paper's regime:
facial recovery of 22.6--25.1\,dB across 5--20\,dB SNR, a correct-minus-wrong token gap of
8.41\,dB (paper 6.46\,dB), and a XuNet steganalyser at chance.

On top of this baseline, the project built \textbf{RACS}, which improves it in two ways:
\begin{itemize}
	\item \textbf{Retrieval instead of generation.} The reference image is found in an indexed
	collection and picked by the token, instead of being painted by a diffusion model. It is
	obtained about $82\times$ faster (0.048\,s instead of 3.92\,s) with no GPU memory, sender and
	receiver get the identical image for all 66 test images, and quality and security are kept or
	slightly improved (PSNR 22.62 vs 22.37\,dB, security gap 8.94 vs 8.67\,dB).
	\item \textbf{Exact recovery.} All 66 secrets are recovered pixel-identical, every wrong key is
	rejected, the chi-square attack is reduced to chance, and Sample Pair Analysis is weakened from
	AUC 1.00 to 0.83. The cover image's realism turned out to matter as much as the embedding rule.
\end{itemize}

A phone prototype shows RACS working on real photographs, with a receive path that needs no AI
model and finishes in one to two seconds.

\section{Learning Outcomes}
\begin{itemize}
	\item Practical understanding of latent diffusion, DDIM, exact inversion with EDICT,
	ControlNet and IP-Adapter, and of where each loses or preserves information.
	\item Deep joint source--channel coding and channel simulation at a controlled SNR.
	\item Multimodal retrieval with CLIP and exact nearest-neighbour search, and why determinism
	matters more than speed when both ends must agree.
	\item Applied cryptography for steganography: key stretching, stream encryption, message
	authentication, and the difference between hiding content and hiding existence.
	\item Steganalysis: chi-square, Sample Pair Analysis and CNN detectors, and how cover-source
	statistics affect detectability.
	\item Research engineering: phase gates, frozen splits, configuration and model hashing, a
	deviation register, and memory management on constrained hardware.
\end{itemize}

\section{Future Work}
\begin{itemize}
	\item \textbf{Full-scale evaluation.} Run the comparison on all 22 images per class and report
	confidence intervals; run EDICT at the paper's 50 steps and strength 0.5.
	\item \textbf{RACS variants.} Vision--language reranking with a fixed deterministic ranking
	(RACS-Agent), lightweight deterministic editing of the retrieved image (RACS-AgentEdit), and a
	retrieve-or-generate fallback when no corpus image is close enough (RACS-Hybrid).
	\item \textbf{Larger corpora.} Grow the corpus from 1{,}000 to 50{,}000--100{,}000 images and
	study the trade-off between retrieval quality and index cost under the determinism contract.
	\item \textbf{Stronger steganalysis.} Evaluate the exact layer against trained CNN
	steganalysers such as SRNet on this cover source, and adapt embedding costs accordingly.
	\item \textbf{Robust exact layer.} Add error-correcting codes so that the exact payload
	survives mild recompression, and support longer keys.
	\item \textbf{Digital branch.} Reproduce the digital semantic codec when a checkpoint becomes
	available.
\end{itemize}

The most natural next step is RACS-Hybrid (Figure~\ref{fig:hybrid}), which keeps retrieval as the
fast default and falls back to diffusion only when the corpus has nothing close enough. The
decision must itself be deterministic and computed from public data, so that the receiver always
takes the same branch as the sender.

\begin{figure}[H]
	\centering
	\begin{tikzpicture}[scale=0.85, transform shape]
		\hyphenpenalty=10000\exhyphenpenalty=10000
		\node[src, text width=2.4cm] (in) at (0,0) {$K_{\mathrm{pub}}$, $X_{\mathrm{feat}}$, token $\tau$};
		\node[newblk, text width=2.6cm] (ret) at (3.6,0) {\textbf{RACS retrieval}\\top-10 pool};
		\node[draw, diamond, aspect=2.2, fill=yellow!10, align=center, font=\small, inner sep=1pt] (dec) at (7.4,0) {best score\\above threshold?};
		\node[newblk, text width=2.4cm] (use) at (11.2,1.2) {\textbf{Use retrieved} $X_{\mathrm{ref}}$\\0.05\,s, no GPU};
		\node[oldblk, text width=2.4cm] (gen) at (11.2,-1.2) {\textbf{Generate} $X_{\mathrm{ref}}$\\diffusion, about 4\,s};
		\draw[arr] (in) -- (ret);
		\draw[arr] (ret) -- (dec);
		\draw[arr] (dec.north) |- node[pos=0.75, above, font=\scriptsize]{yes} (use.west);
		\draw[arr] (dec.south) |- node[pos=0.75, below, font=\scriptsize]{no} (gen.west);
	\end{tikzpicture}
	\caption{Proposed RACS-Hybrid: retrieval by default, generation only as a fallback. The threshold uses public scores only, so both ends make the same choice.}
	\label{fig:hybrid}
\end{figure}

\section{Conclusion}
Diffusion generation is not required to obtain a good steganographic reference image: a
deterministic, indexed retrieval stage delivers the same recovery quality and security at a
small fraction of the cost and without a GPU, provided the transmitter and receiver follow a
strict determinism contract. And approximate generative recovery need not be the end point: a
carefully designed authenticated layer can deliver the exact secret while keeping classical
detectors near chance. Together, these results make coverless semantic steganography cheaper,
exact and more practical for real devices.
